\chapter{Complexity — Path Hard}
%%% generated by the blueprint pipeline from TCSlib.Complexity.SpaceComplexity.Logspace.PathHard
\section{Overview}

Declarations of the Lean module \texttt{TCSlib.Complexity.SpaceComplexity.Logspace.PathHard}.

\section{Declarations}

\begin{definition}[Binary length of a successor]
\lean{Complexity.PathHard.ell}
\label{Complexity.PathHard.ell}
\leanok
For each natural number $m$, define $\ell(m)$ to be the number of bits in the binary
expansion of $m+1$.
\end{definition}

\begin{definition}[Exponential bound]
\lean{Complexity.PathHard.expo}
\label{Complexity.PathHard.expo}
\leanok
\uses{Complexity.PathHard.ell}
For natural numbers $K$, $c_0$, and $m$, define $E(K,c_0,m)=K(c_0+1)\ell(m)+(K+2)$,
where $\ell$ is the length measure used in the surrounding setup.
\end{definition}

\begin{definition}[Number of vertices]
\lean{Complexity.PathHard.nV}
\label{Complexity.PathHard.nV}
\leanok
\uses{Complexity.PathHard.expo}
For natural numbers $K$, $c_0$, and $m$, define $V(K,c_0,m)=2^{E(K,c_0,m)}$, where
$E(K,c_0,m)$ is the bound defined above.
\end{definition}

\begin{lemma}[Bit Length and Binary Size]
\lean{Complexity.PathHard.ell_eq_size}
\label{Complexity.PathHard.ell_eq_size}
\leanok
\uses{Complexity.PathHard.ell}
\difficulty{1}
For every natural number $m$, the length of the binary bit list of $m+1$ equals the
binary size of $m+1$.
\end{lemma}

\begin{lemma}[Logarithmic space bounded by bit length]
\lean{Complexity.PathHard.logSpace_le_ell}
\label{Complexity.PathHard.logSpace_le_ell}
\leanok
\uses{Complexity.PathHard.ell, Complexity.PathHard.ell_eq_size, Complexity.logSpace}
\difficulty{5}
For every natural number $m$, the base-$2$ logarithm of $m$, rounded down, plus $1$ is
at most the number of binary digits of $m+1$.
\end{lemma}

\begin{lemma}[Power Bound from Binary Length]
\lean{Complexity.PathHard.two_pow_ell_le}
\label{Complexity.PathHard.two_pow_ell_le}
\leanok
\uses{Complexity.PathHard.ell, Complexity.PathHard.ell_eq_size}
\difficulty{4}
For every natural number $m$, if $\ell(m)$ is the length of the binary expansion of
$m+1$, then $2^{\ell(m)}\leq 2(m+1)$.
\end{lemma}

\begin{lemma}[Monotonicity of Bit Length]
\lean{Complexity.PathHard.ell_mono}
\label{Complexity.PathHard.ell_mono}
\leanok
\uses{Complexity.PathHard.ell, Complexity.PathHard.ell_eq_size}
\difficulty{3}
For natural numbers $m\le m'$, the length of the binary bit list of $m+1$ is at most the
length of the binary bit list of $m'+1$.
\end{lemma}

\begin{lemma}[Monotonicity of the vertex count]
\lean{Complexity.PathHard.nV_mono}
\label{Complexity.PathHard.nV_mono}
\leanok
\uses{Complexity.PathHard.ell_mono, Complexity.PathHard.expo, Complexity.PathHard.nV}
\difficulty{3}
For nonnegative integers $K,c_0$ and $m\le m'$, the vertex count $n_V(K,c_0,m)$ is at
most $n_V(K,c_0,m')$.
\end{lemma}

\begin{lemma}[Lower Bound on the Vertex Count]
\lean{Complexity.PathHard.four_le_nV}
\label{Complexity.PathHard.four_le_nV}
\leanok
\uses{Complexity.PathHard.expo, Complexity.PathHard.nV}
\difficulty{3}
For all natural numbers $K,c_0,m$, the number $2^{K(c_0+1)\ell(m)+(K+2)}$ is at least
$4$, where $\ell(m)$ is the length of the binary bit list of $m+1$.
\end{lemma}

\begin{lemma}[Upper bound for the vertex count]
\lean{Complexity.PathHard.nV_le}
\label{Complexity.PathHard.nV_le}
\leanok
\uses{Complexity.PathHard.ell, Complexity.PathHard.expo, Complexity.PathHard.nV, Complexity.PathHard.two_pow_ell_le}
\difficulty{4}
For nonnegative integers $K,c_0,m$, let $\ell(m)$ be the length of the binary expansion
of $m+1$, and set $n_V(K,c_0,m)=2^{K(c_0+1)\ell(m)+(K+2)}$. Then \[n_V(K,c_0,m)\le
2^{K+2}\,2^{K(c_0+1)}(m+1)^{K(c_0+1)}.\]
\end{lemma}

\begin{definition}[Machine configuration edges]
\lean{Complexity.PathHard.Edge}
\label{Complexity.PathHard.Edge}
\leanok
\uses{Complexity.Codec.StepRel, Complexity.Codec.acceptLang, Complexity.Codec.initLang, Turing.pairEncode}
Let $N$ be a nondeterministic Turing machine with Boolean alphabet, let $x$ be a finite
Boolean string, and let $s,u,v$ be natural numbers. Declare $(u,v)$ to be an edge at
input $x$ and parameter $s$ if at least one of the following holds: $u=0$, $v\ge2$, and
the encoding of $x$ together with the encoding of the binary representations of $v-2$
and $s$ belongs to the initial language of $N$; or $u\ge2$, $v=1$, and the corresponding
encoding using $u-2$ belongs to its accepting language; or $u,v\ge2$ and $N$ makes a
step from configuration $u-2$ to configuration $v-2$ on input $x$ with parameter $s$.
\end{definition}

\begin{definition}[Encoded adjacency word]
\lean{Complexity.PathHard.adjWordE}
\label{Complexity.PathHard.adjWordE}
\leanok
\uses{Complexity.PathHard.Edge}
Let $N$ be a nondeterministic Turing machine with Boolean alphabet, let $x$ be a finite
Boolean string, and let $s,n$ be natural numbers. Form the Boolean word of length $n^2$
whose entry at each index $k<n^2$ is true exactly when there is an edge at input $x$ and
parameter $s$ from vertex $\lfloor k/n\rfloor$ to vertex $k\bmod n$.
\end{definition}

\begin{lemma}[Adjacency word length]
\lean{Complexity.PathHard.length_adjWordE}
\label{Complexity.PathHard.length_adjWordE}
\leanok
\uses{Complexity.PathHard.adjWordE}
\difficulty{1}
For every finite nondeterministic Turing machine over the Boolean alphabet, every input
bit string, and all natural numbers $s,n$, the adjacency word has length $n^2$.
\end{lemma}

\begin{lemma}[Adjacency word entry]
\lean{Complexity.PathHard.adjWordE_getD}
\label{Complexity.PathHard.adjWordE_getD}
\leanok
\uses{Complexity.PathHard.Edge, Complexity.PathHard.adjWordE}
\difficulty{1}
Let $N$ be a nondeterministic Turing machine over the Boolean alphabet, let $x$ be a
Boolean word, and let $s,n,k$ be natural numbers with $k<n^2$. The entry at position $k$
of the adjacency word for $N,x,s,n$, using false as the default for an out-of-range
position, is true exactly when there is an edge from vertex $\bigl\lfloor
k/n\bigr\rfloor$ to vertex $k \bmod n$ in the graph defined by $N,x,s$.
\end{lemma}

\begin{lemma}[Adjacency word encodes edges]
\lean{Complexity.PathHard.rawEdge_adjWordE}
\label{Complexity.PathHard.rawEdge_adjWordE}
\leanok
\uses{Complexity.PathHard.Edge, Complexity.PathHard.adjWordE, Complexity.PathHard.adjWordE_getD, Complexity.rawEdge}
\difficulty{5}
Let $N$ be a finite nondeterministic Turing machine, let $x$ be a bit string, and let
$s,n,u,v$ be natural numbers. The raw adjacency relation for the word formed by listing
the edge indicators of the graph on $n$ vertices is equivalent to $u<n$, $v<n$, and the
existence of an edge from $u$ to $v$ at space parameter $s$.
\end{lemma}

\begin{definition}[Reduction output]
\lean{Complexity.PathHard.red}
\label{Complexity.PathHard.red}
\leanok
\uses{Complexity.Codec.codeK, Complexity.PathHard.adjWordE, Complexity.PathHard.ell, Complexity.PathHard.nV, Complexity.pathRaw}
Let $N$ be a nondeterministic Turing machine with Boolean alphabet, let $c_0$ be a
natural number, and let $x$ be a finite Boolean string. The output is the path encoding
on $V(K,c_0,|x|)$ vertices, applied to the adjacency word for input $x$, parameter
$c_0\ell(|x|)$, and that same vertex count, with designated vertices $0$ and $1$, where
$K$ is the code parameter of $N$.
\end{definition}

\begin{theorem}[Correctness of the path reduction]
\lean{Complexity.PathHard.red_mem_iff}
\label{Complexity.PathHard.red_mem_iff}
\leanok
\uses{Complexity.Codec.StepRel, Complexity.Codec.acceptLang, Complexity.Codec.codeK, Complexity.Codec.mem_iff_codeChain, Complexity.Codec.mem_of_codePath, Complexity.Codec.pairEncode_inj, Complexity.LogProg.bits_injective, Complexity.PathHard.adjWordE, Complexity.PathHard.ell, Complexity.PathHard.expo, Complexity.PathHard.four_le_nV, Complexity.PathHard.length_adjWordE, Complexity.PathHard.logSpace_le_ell, Complexity.PathHard.nV, Complexity.PathHard.rawEdge_adjWordE, Complexity.PathHard.red, Complexity.logSpace, Complexity.pathRaw, Complexity.pathRawLang, Complexity.rawEdge, Turing.FinNDTM.DecidesInSpace, Turing.pairEncode}
\difficulty{7}
Let $N$ be a finite nondeterministic Turing machine, let $L$ be a language of bit
strings, and let $c_0$ be a natural number. Suppose that on every input $x$, all choice
sequences of some common length $T$ halt, every such computation uses at most
$c_0(\lfloor\log_2(|x|)\rfloor+1)$ work-tape cells, and $N$ accepts $x$ exactly when
$x\in L$. Then, for every bit string $x$, $x\in L$ if and only if the reduction string
constructed from $N,c_0,x$ encodes a path from vertex $0$ to vertex $1$ in the
associated raw graph.
\end{theorem}

\begin{lemma}[Reduction output length]
\lean{Complexity.PathHard.length_red}
\label{Complexity.PathHard.length_red}
\leanok
\uses{Complexity.Codec.codeK, Complexity.PathHard.length_adjWordE, Complexity.PathHard.nV, Complexity.PathHard.red, Complexity.pathRaw, Turing.length_pairEncode}
\difficulty{3}
For every finite nondeterministic Turing machine over the Boolean alphabet, every
natural number $c_0$, and every input bit string $x$, the reduction output has length
exactly $2N^2+2N+7$, where $N$ is the number of vertices used by the reduction on $x$.
\end{lemma}

\begin{lemma}[Polynomial bound on reduction length]
\lean{Complexity.PathHard.length_red_le}
\label{Complexity.PathHard.length_red_le}
\leanok
\uses{Complexity.Codec.codeK, Complexity.PathHard.length_red, Complexity.PathHard.nV, Complexity.PathHard.nV_le, Complexity.PathHard.red}
\difficulty{4}
For every finite nondeterministic Turing machine over the Boolean alphabet and every
natural number $c_0$, there are natural numbers $C,c$ such that, for every input bit
string $x$, the reduction output has length at most $C(|x|+1)^c$.
\end{lemma}

\begin{lemma}[Bits of the reduction string]
\lean{Complexity.PathHard.red_getD}
\label{Complexity.PathHard.red_getD}
\leanok
\uses{Complexity.Codec.codeK, Complexity.PathHard.Edge, Complexity.PathHard.adjWordE_getD, Complexity.PathHard.ell, Complexity.PathHard.length_adjWordE, Complexity.PathHard.nV, Complexity.PathHard.red, Complexity.pairEncode_getD, Complexity.pathRaw, Turing.pairEncode}
\difficulty{5}
Let $N$ be a finite nondeterministic Turing machine, $c_0$ a natural number, $x$ a bit
string, and $i$ a natural number. Write $n$ for the number of vertices in the reduction
and $L=2n$. The bit at position $i$ of the reduction string, taking an out-of-range bit
to be false, equals true for $i<L$, equals false at $i=L$, equals true at $i=L+1$,
equals the edge indicator at row $((i-(L+2))/2)/n$ and column $((i-(L+2))/2)\bmod n$
when $i-(L+2)<2n^2$, and otherwise equals the bit at position $i-(L+2)-2n^2$ of the
five-bit tail $[0,1,0,1,1]$, with out-of-range tail bits false.
\end{lemma}
